\setcounter{section}{0}

\Needspace{5\baselineskip}
\hypertarget{autoregressive-models}{%
\section{Autoregressive Models}\label{autoregressive-models}}

\Needspace{5\baselineskip}
\hypertarget{i.-core-idea}{%
\subsection{I. Core Idea}\label{i.-core-idea}}

The current value in a sequence can be explained and predicted by a linear combination of its values at several preceding time steps plus a random disturbance (error). For example, context can be used to predict the most likely ``next'' word.

\Needspace{5\baselineskip}
A pth-order autoregressive model, denoted AR(p), is defined by:

\[
X_t = c + \phi_1 X_{t-1} + \phi_2 X_{t-2} + \cdots + \phi_p X_{t-p} + \epsilon_t
\]

\Needspace{5\baselineskip}
Interpreting the formula:

\begin{itemize}
\tightlist
\item
  \(X_t\): the variable's value at the current time \(t\) (what we want to predict).
\item
  \(c\): a constant term (intercept).
\item
  \(\phi_1, \phi_2, \ldots, \phi_p\): autoregressive coefficients, the parameters learned by the model. They represent the magnitude and direction of the influence of each past value on the current value.
\item
  \(X_{t-1}, X_{t-2}, \ldots, X_{t-p}\): historical values from the past \(p\) time steps.
\item
  \(\epsilon_t\): a random error term (white noise) at time \(t\), representing fluctuations that historical data cannot explain. Its mean is generally assumed to be 0 and its variance constant.
\end{itemize}

The order p determines how far into the past the model looks. Choosing p is crucial: too small may cause underfitting, while too large may cause overfitting.

Stationarity requirement: the sequence's mean, variance, and covariance should not change substantially over time. A nonstationary sequence generally needs to be made stationary through methods such as differencing (which leads to the ``I'' in ARIMA).

\Needspace{5\baselineskip}
\hypertarget{ii.-applications}{%
\subsection{II. Applications}\label{ii.-applications}}

\begin{enumerate}
\def\labelenumi{\arabic{enumi}.}
\tightlist
\item
  LLMs: treat generation as sequence generation (sequences of words or tokens). The model defines a conditional distribution \(P(x_t | x_{<t})\): given all previously generated words \(x_1, x_2, ..., x_{t-1}\), compute the probability distribution of the next word \(x_t\). Sample from this distribution (or take the most likely value) to obtain \(x_t\). Append the new word \(x_t\) to the context and repeat until an end-of-sequence token is generated or the maximum length is reached.
\end{enumerate}

Implementation: causal masking. In a Transformer, when computing attention at position t, only information from positions 1 through t-1 is visible; ``future'' information is inaccessible, strictly enforcing autoregression.

\begin{enumerate}
\def\labelenumi{\arabic{enumi}.}
\setcounter{enumi}{1}
\tightlist
\item
  Computer vision: very clear, high-quality images can also be generated by ``predicting the next pixel's probability distribution given all previously generated probabilities.'' The disadvantage is extreme slowness: pixels must be generated sequentially, one at a time, without parallelization.
\end{enumerate}

\FloatBarrier
\Needspace{5\baselineskip}
\hypertarget{references-17}{%
\subsection{References}\label{references-17}}

\vspace{0.25em}
\begin{itemize}
\tightlist
\item
  Zhang, A., Lipton, Z. C., Li, M., \& Smola, A. J. (2023). \href{https://D2L.ai}{Dive into Deep Learning}. Cambridge University Press.
\end{itemize}

\clearpage
